\section{Background and Related Work}\label{s:background}

\subsection{LLM-Based Patient Simulation}

We construct this work based on DSM5AgentFlow, which uses LLM agents to generate questionnaire-based assessment dialogues from clinical profiles and produce diagnoses with support from retrieved DSM-5 documents~\cite{ozgun2025trustworthyaipsychotherapymultiagent}.

Other work has used LLMs to simulate patients for clinical and counseling training. Roleplay-doh converts expert feedback into principles that guide simulated mental health patients~\cite{louie2024roleplaydoh}. PATIENT-$\psi$ combines cognitive models with conversational styles for cognitive behavioral therapy training~\cite{wang2024patientpsi}, while EvoPatient develops patient and doctor agents through repeated simulated diagnostic interactions~\cite{du2025evopatient}. These studies show that LLM patient behavior can be shaped through profiles and behavioral instructions.

The present study uses behavioral instructions for a different purpose: it varies the patient's cooperation while keeping the clinical case and questionnaire fixed. This allows the effects of specific resistant response patterns on the dialogue and final diagnosis to be compared under controlled conditions.

\subsection{Clinical Profiles and Grounding}

The patient agent is generated from a clinical profile that specifies a target disorder, relevant symptoms, and background information. The profile provides the reference case for both patient generation and diagnostic evaluation. Patient answers should therefore remain consistent with the profile even when resistance changes how the information is presented.

Profile adherence matters because unsupported symptoms can alter the diagnostic evidence. Roleplay-doh found that initial GPT-4 responses failed to satisfy expert-defined patient principles or dialogue conventions in 20\% of cases~\cite{louie2024roleplaydoh}. In this study, unsupported additions would also make it unclear whether a diagnostic change was caused by resistance or by a change in the simulated case. Section~\ref{s:methodology} describes the prompt boundaries and truth maps used to reduce this problem.

\subsection{Client Resistance}

Psychotherapy research does not use a single definition of client resistance. Behavioral coding frameworks identify resistance from observable moments in which clients limit, divert, or oppose therapeutic work~\cite{chamberlain1984observation,otani1989client}. Other frameworks organize resistance around the therapeutic task being resisted, such as discussing painful affect, recalling personal material, or accepting change~\cite{mahalik1994development}. Motivational interviewing adopts a relational view, treating apparent resistance as discord that can emerge in the interaction rather than as a fixed trait of the client~\cite{miller2013motivational}.

Recent computational research has studied resistance in text-based counseling through both recognition and simulation. PsyFIRE defines 13 resistance types for individual client utterances, and RECAP uses this framework to detect and classify resistance in counseling dialogue~\cite{li2026recap}. ResistClient takes a generative approach: it simulates challenging clients by linking their responses to modeled internal motivations~\cite{liu2026resistclient}. In contrast, this thesis evaluates how 12 predefined resistance subcategories affect information collection and diagnosis in a fixed questionnaire-based assessment.

In this thesis, client resistance refers to observable patient responses that reduce cooperation with the assessment by making the requested information less available, changing how a profile-supported problem is presented, or disrupting the questionnaire process. This definition is closest to the event-level and response-based approaches above, but the categories are reorganized around their effect on diagnostic information collection. The taxonomy is used to control patient-response generation rather than infer a real client's intention or psychological motive. Similar responses in clinical practice may also result from symptoms, anxiety, fatigue, misunderstanding, or difficulty expressing feelings. In this thesis, the resistance label therefore identifies an experimental generation condition rather than a clinical explanation of the patient.

Resistance can also affect outcomes differently depending on how it is expressed. For example, openly hostile resistance has been associated with poorer symptom change and treatment attrition in cognitive behavioral therapy for panic disorder~\cite{schwartz2019resistance}. This supports examining separate resistance categories instead of treating resistance as one uniform condition~\cite{verhulst1990resistance,beutler2002resistance}.

\subsection{Resistance Dimensions and Categories}

Rather than reproducing one existing framework, the taxonomy groups the resistance subcategories according to the part of the questionnaire-based assessment they mainly affect. Content resistance changes the information provided in response to the current item and includes withholding, diverting, and reframing. Interaction resistance interferes more directly with the assessment process and includes contesting and controlling.

Table~\ref{tab:resistance-summary} summarizes the five categories and 12 subcategories. Appendix~\ref{s:appendix-taxonomy} provides the complete definitions, examples, and boundary rules.

\begin{table}[tb]
    \centering
    \small
    \tabcap{Resistance dimensions, categories, and subcategories.}
    \label{tab:resistance-summary}
    \begin{tabularx}{\linewidth}{@{}p{0.18\linewidth} p{0.21\linewidth} X@{}}
        \toprule
        \thead{Dimension} & \thead{Category} & \thead{Subcategories} \\
        \midrule
        Content & Withholding & Vague/minimal information; silence/nonresponse \\
        Content & Diverting & Topic shift/inattention; irrelevant verbosity \\
        Content & Reframing & Minimization; externalization; denial \\
        Interaction & Contesting & Challenging assessment; hostile criticism \\
        Interaction & Controlling & Topic blocking; agenda redirection; demanding diagnosis \\
        \bottomrule
    \end{tabularx}
\end{table}
